
\documentclass[10pt]{article}
\usepackage[utf8]{inputenc}
\usepackage[T1]{fontenc}
\usepackage{amsmath, amssymb, xcolor, geometry, enumitem, multicol, tcolorbox}

\definecolor{mainblue}{RGB}{0, 102, 204}
\geometry{a4paper, margin=1.2cm, top=1cm, bottom=3cm, footskip=1.2cm}

\begin{document}
\enlargethispage{2.5cm}

% Modern Header
\begin{tcolorbox}[colback=mainblue!5, colframe=mainblue, sharp corners, boxrule=0.5pt]
    \small \textbf{Unit:} Unit 0: Foundations of Algebra \hfill \textbf{Name:} \rule{3cm}{0.4pt} \\
    \small \textbf{Topic:} Variables and Expressions / Order of Operations \hfill \textbf{Date:} \rule{2cm}{0.4pt} \\
    \centering \textbf{\large Foundation (Jalapeno)}
\end{tcolorbox}

\vspace{0.2cm}

% MCQ Section
\begin{multicols}{2}
\begin{enumerate}[label=\textbf{Q\arabic*.}, leftmargin=*, itemsep=6pt, parsep=0pt]

  \item \small Evaluate the expression: $2 \cdot 3 + 12 - 8$
  \begin{enumerate}[label=(\Alph*), leftmargin=0.6cm, nosep, itemsep=2pt]
    \item 6
    \item 8
    \item 10
    \item 12
    \item 14
  \end{enumerate}

  \item \small Translate the phrase '5 more than x' into an algebraic expression
  \begin{enumerate}[label=(\Alph*), leftmargin=0.6cm, nosep, itemsep=2pt]
    \item $x - 5$
    \item $x + 5$
    \item $5x$
    \item $5 - x$
    \item $x^2 + 5$
  \end{enumerate}

  \item \small Simplify the expression: $3x + 2x$
  \begin{enumerate}[label=(\Alph*), leftmargin=0.6cm, nosep, itemsep=2pt]
    \item $5x$
    \item $6x$
    \item $x^2$
    \item $x + 5$
    \item $2x - 3$
  \end{enumerate}

  \item \small Evaluate the expression: $(8 - 2) \cdot (3 + 1)$
  \begin{enumerate}[label=(\Alph*), leftmargin=0.6cm, nosep, itemsep=2pt]
    \item 10
    \item 12
    \item 14
    \item 16
    \item 20
  \end{enumerate}

  \item \small A stock price increased by $2.50. If the original price was $x, what is the new price?
  \begin{enumerate}[label=(\Alph*), leftmargin=0.6cm, nosep, itemsep=2pt]
    \item $x - 2.50$
    \item $x + 2.50$
    \item $2.50x$
    \item $x^2 + 2.50$
    \item $x - 2.50x$
  \end{enumerate}

  \item \small A tax of 8% is added to a $500 savings account. What is the total amount after tax?
  \begin{enumerate}[label=(\Alph*), leftmargin=0.6cm, nosep, itemsep=2pt]
    \item $540
    \item $548
    \item $550
    \item $560
    \item $568
  \end{enumerate}

  \item \small Evaluate the expression: $12 - (3 + 2) \cdot 2$
  \begin{enumerate}[label=(\Alph*), leftmargin=0.6cm, nosep, itemsep=2pt]
    \item 0
    \item 2
    \item 4
    \item 6
    \item 8
  \end{enumerate}

  \item \small Translate the phrase '7 less than 3 times a number' into an algebraic expression
  \begin{enumerate}[label=(\Alph*), leftmargin=0.6cm, nosep, itemsep=2pt]
    \item $3x - 7$
    \item $3x + 7$
    \item $7 - 3x$
    \item $7x - 3$
    \item $3x^2 - 7$
  \end{enumerate}

\end{enumerate}
\end{multicols}

\vspace{-0.2cm}
\hrule
\vspace{0.2cm}
\textbf{\small Open Ended Questions (FRQ)}
\begin{enumerate}[label=\textbf{Q\arabic*.}, start=9, itemsep=5pt, topsep=0pt]

  \item \small Evaluate the expression $x^2 + 4x - 5$ when $x = 2$ \vspace{2cm}

  \item \small A company's profit can be modeled by the expression $2x^2 + 5x - 1000$ where $x$ is the number of products sold. Find the profit when $x = 10$ \vspace{2cm}

\end{enumerate}

\vfill

% Chef's Tips Footer
\begin{tcolorbox}[colback=blue!5, colframe=mainblue!50, title=\small \textbf{Chef's Tips}, fonttitle=\bfseries, bottom=1mm]
\begin{itemize}[noitemsep, topsep=2pt, leftmargin=0.5cm]
  \item \scriptsize Tip 1: Remind students to follow the order of operations\item \scriptsize Tip 2: Encourage students to use real-world examples to understand algebraic expressions\item \scriptsize Tip 3: Suggest that students check their work by plugging their answers back into the original equations
\end{itemize}
\end{tcolorbox}

\end{document}
  